\documentclass[12pt,a4paper]{report}

% --- PACKAGES ---
\usepackage[top=1in, bottom=1in, left=1.25in, right=1in]{geometry}
\usepackage{graphicx}
\usepackage{fancyhdr}
\usepackage{setspace}
\usepackage{booktabs}
\usepackage{tabularx}
\usepackage{longtable}
\usepackage{listings}
\usepackage{xcolor}
\usepackage[hidelinks]{hyperref}
\usepackage{amsmath}
\usepackage{caption}
\usepackage{float}
\usepackage{appendix}
\usepackage{titlesec}
\usepackage{parskip}
\usepackage{microtype}
\usepackage{times}
\usepackage{url}
\usepackage{verbatim}
\usepackage{textcomp}

% --- CODE STYLE ---
\definecolor{codegreen}{rgb}{0,0.55,0}
\definecolor{codegray}{rgb}{0.4,0.4,0.4}
\definecolor{codepurple}{rgb}{0.5,0,0.7}
\definecolor{backcolour}{rgb}{0.97,0.97,0.97}
\definecolor{keyblue}{rgb}{0.0,0.2,0.8}

\lstdefinestyle{pythonstyle}{
    backgroundcolor=\color{backcolour},
    commentstyle=\color{codegreen}\itshape,
    keywordstyle=\color{keyblue}\bfseries,
    numberstyle=\tiny\color{codegray},
    stringstyle=\color{codepurple},
    basicstyle=\ttfamily\small,
    breakatwhitespace=false,
    breaklines=true,
    captionpos=b,
    keepspaces=true,
    numbers=left,
    numbersep=6pt,
    showspaces=false,
    showstringspaces=false,
    tabsize=4,
    frame=single,
    rulecolor=\color{gray!40}
}
\lstset{style=pythonstyle, language=Python}

% --- PAGE STYLE ---
\pagestyle{fancy}
\fancyhf{}
\fancyhead[L]{\small\textit{AgriProfit: An ML-Driven Agricultural Price Forecasting and Advisory Platform}}
\fancyhead[R]{\small\textit{Computer Science \& Engineering --- FISAT}}
\fancyfoot[C]{\thepage}
\renewcommand{\headrulewidth}{0.4pt}
\renewcommand{\footrulewidth}{0pt}

% --- GRAPHICS PATH ---
\graphicspath{{report_screenshots/}}

% --- CHAPTER TITLE FORMAT ---
\titleformat{\chapter}[display]
  {\normalfont\large\bfseries\centering}
  {\large Chapter \thechapter}
  {8pt}
  {\Large\bfseries\centering}
\titlespacing*{\chapter}{0pt}{-10pt}{20pt}

% --- LINE SPACING ---
\setstretch{1.5}

% --- RUPEE SYMBOL ---
\newcommand{\rupee}{\textrm{\textit{\textbf{Rs.}}}}

% ============================================================
\begin{document}
% ============================================================

% ============================================================
% TITLE PAGE
% ============================================================
\begin{titlepage}
\thispagestyle{empty}
\begin{center}

\vspace*{0.8cm}
{\LARGE AgriProfit: An ML-Driven Agricultural Price}\\[0.3cm]
{\LARGE Forecasting and Advisory Platform}\\[0.8cm]

{\normalsize\textit{Mini project report submitted in partial fulfilment of the requirements for}}\\[0.1cm]
{\normalsize\textit{the award of the degree of}}\\[1cm]

{\Large\textbf{Bachelor of Technology}}\\[0.2cm]
{\normalsize\textit{in}}\\[0.2cm]
{\Large\textbf{Computer Science \& Engineering}}\\[1.2cm]

{\normalsize Submitted by}\\[0.4cm]
{\normalsize Abhinav K Manoj}\\
{\normalsize Adhwai Shyjith}\\
{\normalsize AdithyaKrishna JG}\\
{\normalsize Al Ameen AK}\\[1.2cm]

\includegraphics[width=5cm]{fisat_logo.jpg}\\[0.3cm]
{\large\textbf{Focus on Excellence}}\\[0.8cm]

{\large\textbf{Federal Institute of Science And Technology (FISAT)\textsuperscript{\textregistered}}}\\
{\normalsize Angamaly, Ernakulam}\\[0.4cm]
{\normalsize\textit{Affiliated to}}\\[0.3cm]
{\large\textbf{APJ Abdul Kalam Technological University}}\\
{\normalsize CET Campus, Thiruvananthapuram}\\[0.3cm]
{\normalsize May 2026}

\end{center}
\end{titlepage}

% ============================================================
% CERTIFICATE PAGE
% ============================================================
\newpage
\thispagestyle{empty}
\begin{center}
{\large\textbf{FEDERAL INSTITUTE OF SCIENCE AND TECHNOLOGY}}\\[0.1cm]
{\large\textbf{(FISAT)}}\\[0.1cm]
{\normalsize Mookkannoor(P.O), Angamaly-683577}\\[0.8cm]
\includegraphics[width=4cm]{fisat_logo.jpg}\\[0.2cm]
{\large\textbf{Focus on Excellence}}\\[1.5cm]
{\Large\textbf{CERTIFICATE}}
\end{center}

\vspace{1.5cm}

\noindent This is to certify that the Mini project report for the project entitled
\textbf{``AgriProfit: An ML-Driven Agricultural Price Forecasting and Advisory Platform''}
is a bonafide report of the project presented during VI\textsuperscript{th} semester
(CSD334 - Mini Project) by \textbf{Abhinav K Manoj (FITCS22002)},
\textbf{Adhwai Shyjith (FITCS22013)}, \textbf{AdithyaKrishna JG (FITCS22017)},
and \textbf{Al Ameen AK (FITCS22018)}, in partial fulfilment of the requirements
for the award of the degree of Bachelor of Technology (B.Tech) in Computer Science
\& Engineering during the academic year 2025--26.

\vspace{3.5cm}

\noindent\hspace{0pt}\makebox[0.33\textwidth][c]{} \makebox[0.33\textwidth][c]{Dr. Neenu Johnson} \makebox[0.33\textwidth][r]{\textbf{Dr. Jyothish K John}}\\[0.5cm]
\noindent\makebox[0.33\textwidth][c]{Project Coordinator} \makebox[0.33\textwidth][c]{Project Guide} \makebox[0.33\textwidth][r]{Head of the Department}

% ============================================================
% ABSTRACT
% ============================================================
\newpage
\thispagestyle{empty}
\chapter*{ABSTRACT}
\addcontentsline{toc}{chapter}{ABSTRACT}

AgriProfit is an ML-driven agricultural price forecasting and advisory platform designed to
empower Indian farmers with data-driven market intelligence. The platform addresses the
fundamental information asymmetry in Indian agricultural commodity markets, where large
traders hold significant informational advantages over smallholder farmers who produce the
bulk of the country's food.

The core of the system is a 7-day ahead price forecasting engine built on direct multi-step
XGBoost models trained on approximately 25 million price records from the AGMARKNET
government database spanning 10 years (2014--2024). The model employs 18 carefully
engineered features including temporal lag values (1, 2, 3, 7, 14, 21, and 30 days),
rolling statistics over 7, 14, and 30-day windows, and calendar features capturing
weekly and seasonal cycles. A log1p price transformation ensures stable training across
commodities with vastly different price scales. Seven separate XGBoost regressors are
trained per commodity --- one per forecast horizon --- achieving a median 7th-horizon
MAPE of 14.4\% across 60 successfully trained commodities.

The platform integrates a Harvest Advisor module that recommends optimal crops for a
farmer's district and current season. It uses Random Forest yield prediction models
trained on 232,858 rows of real government yield data covering 102 crops across 646
districts from 1997--2015, combined with market price forecasts, to compute expected
profit per hectare. Weather warnings from the Open-Meteo API are integrated to flag
agricultural risks.

A Transport Logistics Engine helps farmers choose the most profitable mandi by computing
diesel-adjusted freight costs using real road distances (OSRM routing API), commodity-specific
exponential spoilage models, regional hamali (loading/unloading) charges, and market risk
scores. The platform also provides inventory tracking, sales analytics, community forums,
and arbitrage detection between mandis.

Built with FastAPI (Python 3.11), Next.js 14 (TypeScript), and PostgreSQL 15, the platform
exposes 18 RESTful API route modules. The system achieves a 100\% test pass rate across
598 automated tests with 61.37\% statement coverage. The platform directly supports
SDG~1 (No Poverty) and SDG~2 (Zero Hunger) by improving farmer incomes through
better market decision-making.

% ============================================================
% CONTRIBUTION BY AUTHOR
% ============================================================
\newpage
\thispagestyle{empty}
\chapter*{Contribution by Author}
\addcontentsline{toc}{chapter}{Contribution by Author}

\vspace{3cm}

\noindent\textbf{Abhinav K Manoj} designed and implemented the core ML forecasting pipeline,
including the 7-day direct multi-step XGBoost model architecture, feature engineering module
(\texttt{features\_7d.py}), model training scripts, and the forecast serving layer with
empirical per-horizon confidence bands derived from holdout residuals.

\vspace{0.6cm}
\noindent\textbf{Adhwai Shyjith} developed the Next.js 14 frontend web application,
including the interactive dashboard, price forecast visualizations with Recharts,
harvest advisor UI, community forum interface, and mobile-responsive design system
using Tailwind CSS.

\vspace{0.6cm}
\noindent\textbf{AdithyaKrishna JG} built the FastAPI backend REST API, PostgreSQL
database schema with Alembic migrations, SQLAlchemy ORM models, JWT-based OTP
authentication system, and the inventory and sales tracking modules.

\vspace{0.6cm}
\noindent\textbf{Al Ameen AK} implemented the transport logistics engine (freight
calculation, spoilage modelling, risk scoring), the AGMARKNET data integration and
synchronisation pipeline, the harvest advisor backend module, and the OSRM-based
road distance routing service with database caching.

\vspace{4cm}

\begin{flushright}
Abhinav K Manoj\\
Adhwai Shyjith\\
AdithyaKrishna JG\\
Al Ameen AK
\end{flushright}

% ============================================================
% ACKNOWLEDGMENT
% ============================================================
\newpage
\thispagestyle{empty}
\chapter*{ACKNOWLEDGMENT}
\addcontentsline{toc}{chapter}{ACKNOWLEDGMENT}

\vspace{2cm}

We express our sincere gratitude to \textbf{Dr.~Neenu Johnson}, our Project Guide, for
her invaluable guidance, continuous encouragement, and expert technical advice throughout
the development of this project. Her insights into machine learning methodologies and
agricultural data systems significantly shaped the direction of our work.

We extend our heartfelt thanks to \textbf{Dr.~Jyothish K John}, Head of the Department
of Computer Science and Engineering, for providing an environment conducive to innovation
and for his support in facilitating this mini project.

We are grateful to the management of \textbf{Federal Institute of Science and Technology
(FISAT)}, Angamaly, for providing us with the necessary infrastructure, computing resources,
and academic support throughout the course of this programme.

We also acknowledge the \textbf{Government of India's AGMARKNET} portal (data.gov.in)
and the National Horticulture Board for providing open-access agricultural commodity price
data and crop yield records that form the backbone of this project's datasets.

Finally, we thank our families and colleagues for their patience and moral support during
the course of this project.

\begin{flushright}
Abhinav K Manoj\\
Adhwai Shyjith\\
AdithyaKrishna JG\\
Al Ameen AK
\end{flushright}

% ============================================================
% TOC / LOF / LOT
% ============================================================
\newpage
\tableofcontents
\newpage
\listoffigures
\newpage
\listoftables

% ============================================================
% CHAPTER 1: INTRODUCTION
% ============================================================
\chapter{Introduction}

\section{Overview}

Agriculture is the backbone of the Indian economy, employing approximately 54\% of the
workforce and contributing nearly 17\% to the national GDP~\cite{agricultureGDP}. Despite
its enormous scale, the agricultural sector faces a chronic and widening information
asymmetry: large commodity traders and market intermediaries hold significant informational
advantages, while the 100 million smallholder farmers who produce the bulk of India's food
operate largely without reliable market intelligence.

The consequences are severe and well-documented. Indian farmers receive on average only
20--30\% of the consumer price for their produce, with the remainder captured by supply
chain intermediaries who exploit informational advantages~\cite{farmersIncome}. A farmer
who planted onions in October without price guidance may face a 60\% price drop by December,
while a farmer who sells wheat today without knowing that prices will rise by 15\% next week
leaves significant income on the table.

AgriProfit was conceived to bridge this information gap through a comprehensive digital
platform that applies modern machine learning to publicly available government market data.
The platform is built on three principal observations:

\begin{enumerate}
    \item \textbf{Data exists but is inaccessible}: The AGMARKNET government portal contains
    10+ years of daily price data for 256+ commodities across 3,000+ mandis in India, but
    this data is not presented in a form that enables price forecasting or decision support
    for individual farmers.
    \item \textbf{ML can provide reliable short-term forecasts}: Agricultural commodity prices
    exhibit learnable patterns --- seasonal cycles, weekly rhythms, and local supply-demand
    dynamics --- that gradient boosting models can capture effectively.
    \item \textbf{Multiple decisions need simultaneous support}: Farmers need guidance not
    just on current prices, but on which crop to plant, when to harvest, and which mandi
    offers the best net returns after transport costs.
\end{enumerate}

The platform provides the following core capabilities:

\begin{itemize}
    \item \textbf{7-Day Price Forecasting}: Direct multi-step XGBoost models provide
    day-by-day price predictions for 63+ commodities in the farmer's district, with
    empirical 80\% confidence bands.
    \item \textbf{Harvest Advisor}: Recommends the top-5 crops for a farmer's district
    and current month, based on predicted yields (Random Forest) combined with forecasted
    prices to compute expected profit per hectare.
    \item \textbf{Transport Optimizer}: Compares net profit from selling at different
    mandis after accounting for vehicle freight costs, commodity spoilage during transport,
    hamali charges, and market risk scores.
    \item \textbf{Real-time Price Tracking}: Automated daily synchronisation from the
    AGMARKNET API keeps the price database current.
    \item \textbf{Community Platform}: A community forum enables farmers to share local
    market knowledge, pest or disease alerts, and farming techniques.
    \item \textbf{Inventory and Sales Analytics}: Farmers can track produce holdings and
    analyse their sales history to identify optimal selling patterns.
    \item \textbf{Arbitrage Detection}: Identifies price differentials between mandis
    exceeding transport costs, flagging profitable arbitrage opportunities.
\end{itemize}

The system is built with a Python FastAPI backend, a Next.js TypeScript frontend, and a
PostgreSQL database, structured as a modern web application accessible from both desktop
and mobile browsers. Ethical considerations include presenting forecasts with explicit
uncertainty bands to prevent over-reliance, using only publicly available government data,
and implementing rate limiting and OTP-based authentication to protect farmer privacy.

\begin{figure}[H]
    \centering
    \includegraphics[width=0.95\textwidth]{ss_dashboard.png}
    \caption{AgriProfit Main Dashboard}
    \label{fig:dashboard}
\end{figure}

\section{Problem Statement}

The Indian agricultural market suffers from a fundamental information asymmetry problem
that disproportionately harms smallholder farmers.

\textbf{Background of the Problem}: Agricultural commodity prices in India are highly
volatile, driven by monsoon variability, transportation bottlenecks, supply-demand
imbalances, and government policy interventions. The AGMARKNET database records daily
prices across thousands of mandis, yet this data is not transformed into actionable
forecasts for farmers. Existing solutions such as the Kisan Suvidha app and eNAM
provide historical prices but offer no forward-looking predictions. Studies show that
Indian farmers receive only 20--30\% of consumer prices for their produce due to
information asymmetry and intermediary exploitation~\cite{farmersIncome}.

\textbf{Anchor}: The resolution requires building an ML pipeline that ingests historical
price data, learns commodity-specific price dynamics, and generates 7-day ahead predictions
accurate enough (MAPE $<$ 20\%) to support real farming decisions --- combined with crop
advisory and transport optimisation tools to address the complete decision lifecycle.

\textbf{General Problem}: At the national scale, price uncertainty causes Indian farmers
to collectively lose an estimated \rupee~50,000--80,000 crore annually through suboptimal
selling decisions, distress sales, and uninformed crop selection.

\textbf{Specific Problem}: In Kerala and Karnataka agricultural districts --- the primary
focus of this platform --- commodity price fluctuations of 40--60\% within a single season
are common. Farmers lack digital tools to anticipate price movements, compare mandi options
after transport costs, or identify crops to plant based on expected market conditions at
harvest time.

\section{Objectives}

The objectives of this project are:
\begin{enumerate}
    \item Develop a 7-day ahead price forecasting system achieving median MAPE $\leq$ 20\%
    for trained commodities, using direct multi-step XGBoost models with engineered
    lag and calendar features.
    \item Build a Harvest Advisor module that recommends crops based on district suitability,
    seasonal crop calendar, Random Forest yield predictions, and forecasted market prices.
    \item Implement a Transport Logistics Engine that identifies the most profitable mandi
    for a farmer's commodity, accounting for freight cost, spoilage, hamali, and market risk.
    \item Create a real-time AGMARKNET data synchronisation pipeline that maintains a
    current price history database with automated daily updates.
    \item Design an intuitive, responsive web application accessible to users with basic
    smartphone literacy, available in both desktop and mobile layouts.
    \item Implement automated testing with $\geq$ 60\% statement coverage and 100\% pass
    rate to ensure system reliability before deployment.
\end{enumerate}

\section{Scope of the Project}

The scope of AgriProfit encompasses:
\begin{itemize}
    \item \textbf{Geographic Scope}: Indian agricultural districts with price data available
    in AGMARKNET, with focus on Kerala and Karnataka.
    \item \textbf{Commodity Scope}: 63+ commodity types trained with ML models, covering
    vegetables, fruits, pulses, cereals, oilseeds, spices, and flowers.
    \item \textbf{Forecasting Scope}: 7-day ahead price predictions with per-horizon
    empirical confidence bands (80\% coverage).
    \item \textbf{User Scope}: Farmers, agricultural traders, cooperatives, and government
    agricultural extension officers.
    \item \textbf{Data Scope}: 10 years of AGMARKNET price history (2014--2024) and
    government crop yield records (1997--2015).
\end{itemize}

Out of scope: satellite imagery analysis, IoT sensor integration, real-time satellite
weather forecasting, and agri-loan or insurance processing.

\section{Social Relevance}

AgriProfit directly supports two United Nations Sustainable Development Goals:
\begin{itemize}
    \item \textbf{SDG 1 --- No Poverty}: By improving farmers' ability to sell at optimal
    prices and reduce information-disadvantaged transactions, the platform can increase
    smallholder income by 10--20\%.
    \item \textbf{SDG 2 --- Zero Hunger}: By guiding farmers toward high-value crops with
    good market prospects, the system supports both farmer livelihoods and efficient food
    production allocation.
\end{itemize}

The platform also supports \textbf{SDG 8} (Decent Work and Economic Growth) through digital
empowerment of the agricultural workforce, and \textbf{SDG 9} (Industry, Innovation and
Infrastructure) by demonstrating applied ML in the agricultural sector.

\section{Organization of the Report}

This mini project report is constructed with 7 chapters. Chapter~1 explains the problem
statement, objectives, and scope of the project. Chapter~2 presents a literature review
on agricultural price forecasting, machine learning for time series, and related digital
agriculture platforms. Software requirement specifications, system architecture, database
design, and algorithms are explained in Chapter~3. Chapter~4 gives implementation details
including the ML pipeline, backend API, frontend application, and deployment. Testing
strategies and test cases are explained in Chapter~5. Chapter~6 discusses the forecast
accuracy results, comparison with existing systems, and sustainability impact assessment.
Finally, Chapter~7 gives the conclusion and future scope of the project work.

% ============================================================
% CHAPTER 2: LITERATURE REVIEW
% ============================================================
\chapter{Literature Review}

Machine learning for agricultural price forecasting and digital advisory platforms for
farmers has been an active research area. This chapter reviews the most relevant prior work
and positions AgriProfit's contributions.

\section{Related Work}

\textbf{Temporal Pattern Learning in Commodity Markets}: Soni and Sharma~\cite{lstm2020}
proposed LSTM-based recurrent neural networks for multi-step agricultural commodity price
forecasting in Indian markets. Their model achieved competitive MAPE on some commodities,
but LSTM approaches require significantly more training data per commodity and are prone
to error accumulation in recursive multi-step prediction. Their work demonstrated that
lag-based features --- specifically 7-day and 30-day lags --- are among the most predictive
for agricultural prices.

\textbf{Gradient Boosting for Time Series}: Chen and Guestrin~\cite{xgboostseries}
introduced XGBoost and demonstrated that gradient-boosted trees with lag features
outperform ARIMA and LSTM on tabular time-series problems when the training dataset
per series is limited. Their direct multi-step approach --- training separate models for
each forecast horizon --- avoids recursive error accumulation and provides better-calibrated
forecasts for horizons beyond three days.

\textbf{Prophet for Agricultural Forecasting}: Taylor and Letham~\cite{taylor2018forecasting}
introduced Facebook Prophet, a decomposable time-series model well-suited to series with
strong seasonal patterns. Singh and Srivastava~\cite{prophetagri} applied Prophet to Indian
vegetable price forecasting, showing strong performance on commodities with clear annual and
weekly seasonality (onion, tomato). However, Prophet requires national-level aggregation
for reliable decomposition, making it less suitable for district-level forecasting.

\textbf{Digital Platforms for Agricultural Advisory}: Sinha et al.~\cite{digitalagri2021}
reviewed 47 digital agriculture platforms deployed in South and Southeast Asia. They found
that platforms combining price information with actionable recommendations (crop choice,
selling timing) achieved 3$\times$ higher farmer adoption rates than price-only platforms.
Their review identifies transport cost integration as a key missing feature in most existing
systems.

\textbf{Yield Prediction using Machine Learning}: Devi and Venkatesan~\cite{yieldrf2019}
demonstrated that Random Forest models trained on district-level yield data achieve
competitive accuracy with deep learning approaches while requiring significantly less
training data and being more interpretable for extension officers.

\section{Comparison of Related Works}

\begin{table}[H]
\centering
\caption{Comparison of Related Work with AgriProfit}
\label{tab:litcomp}
\begin{tabular}{|l|c|c|c|c|c|}
\hline
\textbf{System} & \textbf{Forecast} & \textbf{Crop Adv.} & \textbf{Transport} & \textbf{Community} & \textbf{MAPE} \\
\hline
LSTM \cite{lstm2020} & 7-day & No & No & No & 18.2\% \\
Prophet \cite{prophetagri} & 30-day & No & No & No & 21.5\% \\
eNAM & Historical & No & No & No & N/A \\
Kisan Suvidha & Historical & Partial & No & No & N/A \\
\textbf{AgriProfit (ours)} & \textbf{7-day} & \textbf{Yes} & \textbf{Yes} & \textbf{Yes} & \textbf{14.4\%} \\
\hline
\end{tabular}
\end{table}

\section{Proposed System}

Based on gaps identified in the literature, AgriProfit proposes the following novel
contributions:

\begin{enumerate}
    \item \textbf{Direct Multi-Step XGBoost Forecasting}: Unlike recursive approaches,
    we train 7 separate XGBoost regressors per commodity (one per horizon), eliminating
    error accumulation and enabling per-horizon uncertainty quantification using empirical
    residual quantiles from the holdout set.
    \item \textbf{Integrated Transport-Aware Advisory}: We combine price forecasting with
    real-time road distance calculation (OSRM), diesel-adjusted freight economics, and
    commodity-specific spoilage modelling to provide net-profit-optimal mandi recommendations.
    \item \textbf{Full-Stack Production Platform}: Unlike research systems, AgriProfit
    provides a production-ready web application with OTP authentication, persistent
    inventory tracking, community features, and automated daily data synchronisation.
    \item \textbf{Harvest-to-Sale Pipeline}: The system covers the complete farmer decision
    timeline from crop selection (months before planting) through price monitoring (during
    growing season) to optimal mandi selection (at harvest time).
\end{enumerate}

% ============================================================
% CHAPTER 3: DESIGN METHODOLOGIES
% ============================================================
\chapter{Design Methodologies}

\section{Software Requirement Specifications}

A software requirements specification is a detailed description of the intended purpose,
features, functionality, and other aspects of a software application. It serves as a guide
and foundational reference for the development team.

\subsection{Functional Requirements}

\begin{enumerate}
    \item \textbf{User Authentication}: The system shall support OTP-based phone number
    authentication. Users shall be assigned roles (farmer, trader, admin). JWT tokens
    shall expire after a configurable period.
    \item \textbf{Price Forecasting}: The system shall provide 7-day ahead modal price
    forecasts for at least 60 commodity types in the user's selected district, with
    low/high confidence bands representing 80\% empirical coverage.
    \item \textbf{Harvest Advisory}: The system shall recommend the top-5 crops for a
    given district and current month, displaying predicted yield, expected price, and
    profit per hectare.
    \item \textbf{Transport Comparison}: Given a commodity, quantity, and origin district,
    the system shall return a ranked list of mandis by estimated net profit after freight
    and spoilage deduction.
    \item \textbf{Price History}: The system shall display historical price charts for any
    commodity-district combination for the past 90 days.
    \item \textbf{Inventory Management}: Users shall be able to log commodity holdings with
    quantity, purchase price, and purchase date.
    \item \textbf{Community Forum}: Users shall be able to create, read, and reply to posts
    about agricultural topics.
    \item \textbf{Notifications}: The system shall send in-app alerts when prices move more
    than 10\% above or below forecast.
    \item \textbf{Admin Panel}: Administrators shall view system statistics and trigger
    manual data synchronisation.
    \item \textbf{Data Synchronisation}: The system shall automatically synchronise price
    data from AGMARKNET API on a configurable schedule.
\end{enumerate}

\subsection{Non-Functional Requirements}

\begin{enumerate}
    \item \textbf{Performance}: The forecast API endpoint shall respond within 500~ms for
    cached requests and within 3~seconds for uncached requests.
    \item \textbf{Scalability}: The API layer shall support horizontal scaling via
    containerisation (Docker).
    \item \textbf{Security}: All endpoints require JWT authentication. OTP login implements
    rate limiting (5 attempts per hour). SQL injection is prevented through SQLAlchemy
    parameterised queries.
    \item \textbf{Usability}: The web interface shall be responsive and functional on
    screen widths from 320~px (mobile) to 1920~px (desktop).
    \item \textbf{Reliability}: Forecast failures shall fall back gracefully to seasonal
    statistics rather than returning server errors.
    \item \textbf{Maintainability}: Code coverage shall be maintained at $\geq$ 60\% with
    automated tests. All API endpoints shall be documented via auto-generated OpenAPI/Swagger.
    \item \textbf{Data Integrity}: The price history table enforces a unique constraint on
    (commodity\_id, mandi\_name, price\_date) to prevent duplicate records.
\end{enumerate}

\section{Software Design Document}

\subsection{System Features}

The system provides the following major features:
\begin{itemize}
    \item \textbf{ML Forecasting Engine}: XGBoost-based 7-day price prediction with confidence bands
    \item \textbf{Harvest Advisor}: Crop recommendation combining yield prediction and price forecasting
    \item \textbf{Transport Logistics Engine}: Multi-mandi profit comparison
    \item \textbf{AGMARKNET Integration}: Automated price data synchronisation
    \item \textbf{User Management}: OTP auth, role-based access, profile management
    \item \textbf{Community Platform}: Posts, replies, and moderation
    \item \textbf{Inventory \& Sales Tracking}: Commodity holdings and sales analytics
    \item \textbf{Arbitrage Detection}: Cross-mandi price differential identification
    \item \textbf{Admin Dashboard}: System monitoring and management
\end{itemize}

\subsection{System Architecture Design}

AgriProfit follows a three-tier client-server architecture:

\begin{itemize}
    \item \textbf{Presentation Tier}: Next.js 14 application with React components,
    Tailwind CSS, and Recharts for data visualisation. Communicates with the API tier
    via an Axios HTTP client.
    \item \textbf{Application Tier}: FastAPI Python application exposing 18 RESTful
    route modules. Handles business logic, ML model serving, authentication, and
    scheduled data sync via APScheduler.
    \item \textbf{Data Tier}: PostgreSQL 15 relational database with approximately
    25 million price records. ML model artefacts stored as \texttt{.joblib} files.
    Parquet history files used as fallback data source for thin-district scenarios.
\end{itemize}

External integrations:
\begin{itemize}
    \item AGMARKNET API (data.gov.in) for daily price data
    \item OSRM (Open Source Routing Machine) for road distance calculation
    \item Open-Meteo API for weather forecast integration
\end{itemize}

\subsection{Constraints}

\begin{itemize}
    \item \textbf{Technology}: ML models require Python 3.11+ and are not
    browser-executable; a server-side runtime is mandatory.
    \item \textbf{Data}: AGMARKNET API has rate limits and occasionally provides
    incomplete data for small mandis; robust fallback logic is required.
    \item \textbf{Hardware}: Model training requires $\geq$ 16~GB RAM for the
    25M-row dataset. Model serving requires $\geq$ 4~GB RAM to hold 63 XGBoost models
    in memory.
    \item \textbf{Network}: OSRM road distance queries take 1--2 seconds each;
    a database caching layer prevents repeated computation.
\end{itemize}

\subsection{Application Architecture Design}

The backend follows a layered architecture within each feature module:

\begin{itemize}
    \item \textbf{Routes Layer} (\texttt{routes.py}): FastAPI route definitions,
    Pydantic request validation, rate limiting via slowapi.
    \item \textbf{Service Layer} (\texttt{service.py}): Business logic, ML model
    invocation, external API calls, caching.
    \item \textbf{Schema Layer} (\texttt{schemas.py}): Pydantic request/response
    models for API contract enforcement.
    \item \textbf{Model Layer} (\texttt{models/}): SQLAlchemy ORM table definitions
    with typed mapped\_column declarations.
\end{itemize}

Each feature (forecast, harvest\_advisor, transport, inventory, community, etc.) is
a self-contained module following this structure, enabling independent development
and testing.

\subsection{API Design}

An API (Application Programming Interface) is a set of rules enabling software
applications to communicate with each other. Table~\ref{tab:api} lists the primary
REST API endpoints.

\begin{table}[H]
\centering
\caption{Key REST API Endpoints}
\label{tab:api}
\begin{tabular}{|l|l|p{5.5cm}|}
\hline
\textbf{Method} & \textbf{Endpoint} & \textbf{Description} \\
\hline
POST & /api/v1/auth/request-otp & Request OTP for phone-based login \\
POST & /api/v1/auth/verify-otp & Verify OTP and receive JWT tokens \\
GET  & /api/v1/forecast/\{commodity\}/\{district\} & 7-day price forecast \\
GET  & /api/v1/commodities & List all commodities \\
GET  & /api/v1/mandis & List mandis with filters \\
GET  & /api/v1/prices/latest & Latest commodity prices by district \\
POST & /api/v1/transport/compare & Compare mandi net profits \\
GET  & /api/v1/harvest-advisor/\{district\} & Crop recommendations \\
GET  & /api/v1/inventory & User's inventory items \\
POST & /api/v1/inventory & Add inventory item \\
POST & /api/v1/community/posts & Create community post \\
GET  & /api/v1/analytics/price-trend & Price trend analytics \\
GET  & /api/v1/arbitrage & Detect mandi arbitrage \\
GET  & /api/v1/admin/stats & System statistics (admin only) \\
\hline
\end{tabular}
\end{table}

\begin{figure}[H]
    \centering
    \includegraphics[width=\textwidth]{08_swagger.png}
    \caption{Auto-generated Swagger API Documentation}
    \label{fig:swagger}
\end{figure}

\subsection{Database Design}

The PostgreSQL database contains the following primary tables:

\begin{table}[H]
\centering
\caption{Database Schema Summary}
\label{tab:schema}
\begin{tabular}{|l|p{8.5cm}|}
\hline
\textbf{Table} & \textbf{Key Columns} \\
\hline
\texttt{users} & id (UUID PK), phone\_number, role, name, district, state \\
\texttt{commodities} & id (UUID PK), name, category, slug, entity\_id \\
\texttt{mandis} & id (UUID PK), name, state, district, market\_code (UNIQUE) \\
\texttt{price\_history} & id (UUID PK), commodity\_id (FK), mandi\_id (FK), price\_date, modal\_price, min\_price, max\_price \\
\texttt{forecast\_cache} & id (UUID PK), commodity, district, forecast\_date, predictions (JSONB) \\
\texttt{inventory} & id (UUID PK), user\_id (FK), commodity\_id (FK), quantity, purchase\_price \\
\texttt{sale} & id (UUID PK), user\_id (FK), commodity\_id (FK), quantity, sale\_price, mandi\_id \\
\texttt{crop\_yields} & id (UUID PK), district, state, crop, year, yield\_kg\_ha, area\_ha \\
\texttt{community\_post} & id (UUID PK), user\_id (FK), title, body, created\_at \\
\texttt{road\_distance\_cache} & id (UUID PK), origin, destination, distance\_km, source \\
\hline
\end{tabular}
\end{table}

The \texttt{price\_history} table is the largest with $\approx$25 million rows. All
queries on this table include date-range filters. The \texttt{forecast\_cache} table uses
JSONB for flexible storage of per-horizon predictions with confidence bands.

\subsection{Technology Stack}

\begin{table}[H]
\centering
\caption{Technology Stack}
\label{tab:techstack}
\begin{tabular}{|l|l|l|}
\hline
\textbf{Layer} & \textbf{Technology} & \textbf{Version} \\
\hline
Frontend Framework & Next.js & 14.x \\
Frontend Language & TypeScript & 5.x \\
UI Styling & Tailwind CSS & 3.x \\
Charts & Recharts & 2.x \\
Backend Framework & FastAPI & 0.115.x \\
Backend Language & Python & 3.11 \\
ORM & SQLAlchemy & 2.0.x \\
Database & PostgreSQL & 15 \\
ML --- Forecasting & XGBoost & 2.1.x \\
ML --- Seasonal Fallback & Prophet & 1.1.x \\
ML --- Yield Prediction & scikit-learn (Random Forest) & 1.5.x \\
Data Processing & pandas / NumPy & 2.x / 1.26.x \\
Authentication & PyJWT & 2.x \\
Scheduler & APScheduler & 3.x \\
Rate Limiting & slowapi & 0.1.x \\
Backend Testing & pytest & 8.x \\
Frontend Testing & Vitest + React Testing Library & 1.x \\
\hline
\end{tabular}
\end{table}

\section{Use Case Diagrams / Data Flow Diagrams}

\textbf{Primary Use Cases:}

\begin{itemize}
    \item \textbf{UC-01 View Price Forecast}: Farmer selects commodity and district
    $\rightarrow$ System returns 7-day forecast chart with confidence bands.
    \item \textbf{UC-02 Get Crop Recommendation}: Farmer selects district $\rightarrow$
    System returns top-5 crops ranked by expected profit per hectare.
    \item \textbf{UC-03 Compare Mandis}: Farmer enters commodity, quantity, and origin
    $\rightarrow$ System returns ranked mandi list with net profit estimates.
    \item \textbf{UC-04 Track Inventory}: Farmer logs purchased produce $\rightarrow$
    System records inventory and links to price history for P\&L calculation.
    \item \textbf{UC-05 Community Post}: Farmer posts market observation $\rightarrow$
    Other farmers in same district receive notification.
\end{itemize}

\textbf{Data Flow for Price Forecast (UC-01):}
\begin{enumerate}
    \item Frontend sends \texttt{GET /api/v1/forecast/\{commodity\}/\{district\}}
    \item ForecastService7D checks \texttt{forecast\_cache} table (6-hour TTL)
    \item On cache miss: loads XGBoost v5 model and metadata from disk
    \item Queries PostgreSQL \texttt{price\_history} for last 60 days of district prices
    \item If fewer than 14 days available, falls back to Parquet history artefact
    \item Builds 18-feature serving vector using \texttt{build\_serving\_vector()}
    \item Runs 7 separate XGBoost \texttt{predict()} calls (horizons h1 to h7)
    \item Applies empirical residual quantiles for p10/p90 price bands
    \item Stores result in \texttt{forecast\_cache}; returns \texttt{ForecastResponse} JSON
\end{enumerate}

\section{Algorithms}

\textbf{Algorithm~1: Direct Multi-Step XGBoost Price Forecasting}

\begin{lstlisting}[caption={Feature Vector Construction for 7-Day Forecast (features\_7d.py)}]
LAG_DAYS = [1, 2, 3, 7, 14, 21, 30]
ROLL_WINDOWS = [7, 14, 30]

def build_serving_vector(series, district_enc, target_date):
    df = pd.DataFrame({"price": series.values},
                      index=pd.to_datetime(series.index))
    # Log-transform for scale invariance
    df["log_price"] = np.log1p(df["price"].clip(lower=0.0))

    # Lag features
    for d in LAG_DAYS:
        df[f"lag_{d}"] = df["log_price"].shift(d)

    # Rolling statistics (shift 1 to prevent leakage)
    rolled = df["log_price"].shift(1)
    for w in ROLL_WINDOWS:
        df[f"roll_mean_{w}"] = rolled.rolling(w).mean()
        df[f"roll_std_{w}"] = rolled.rolling(w).std().fillna(0)

    # Calendar features
    df["day_of_week"]  = df.index.dayofweek
    df["month"]        = df.index.month
    df["day_of_year"]  = df.index.day_of_year
    df["week_of_year"] = df.index.isocalendar().week
    df["district_enc"] = district_enc

    return df[FEATURE_COLS].iloc[[-1]]  # last row = today

# Inference: 7 models, one per horizon
predictions = []
for h in range(1, 8):
    log_pred   = models[h].predict(feature_vector)
    price_pred = np.expm1(log_pred)   # inverse log1p
    predictions.append(float(price_pred[0]))
\end{lstlisting}

\textbf{Algorithm~2: Harvest Advisor Crop Recommendation}

\begin{enumerate}
    \item Filter \texttt{CROP\_CALENDAR} for crops suitable for the current month in
    the given district.
    \item For each eligible crop:
    \begin{enumerate}
        \item Load the appropriate Random Forest yield model (by crop category)
        \item Predict yield (kg/ha) from district soil and climate features
        \item Invoke \texttt{ForecastService} to get expected price at harvest month
        \item Compute:\quad
        $\texttt{profit\_per\_ha} = \dfrac{\texttt{yield\_kg\_ha}}{100} \times
        \texttt{price\_per\_quintal} - \texttt{input\_cost}$
    \end{enumerate}
    \item Rank crops by expected profit per hectare.
    \item Attach weather warnings from Open-Meteo API.
    \item Return top-5 recommendations with full explanation.
\end{enumerate}

\textbf{Algorithm~3: Transport Mandi Comparison}

\begin{enumerate}
    \item For each candidate mandi within the configured radius:
    \begin{enumerate}
        \item Compute road distance: DB cache $\rightarrow$ OSRM API $\rightarrow$
        Haversine $\times$ 1.35 fallback
        \item Select optimal vehicle type using 90\% practical capacity rule
        \item Compute freight: base rate $\times$ distance + toll + diesel surcharge +
        breakdown reserve (5\%)
        \item Compute spoilage: $\texttt{loss} = Q \times (1-(1-r)^{h/24})$,
        where $Q$ = quantity, $r$ = daily spoilage rate, $h$ = travel hours
        \item Compute hamali charges by region (North/South/Maharashtra/default)
        \item Fetch latest modal price for commodity at this mandi
        \item Compute net profit: $\texttt{gross} - \texttt{freight} -
        \texttt{spoilage} - \texttt{hamali}$
        \item Compute risk score from price volatility, data staleness, and sample size
    \end{enumerate}
    \item Sort mandis by net profit; apply guardrails to flag unsafe recommendations.
    \item Return ranked list with full cost breakdown and structured audit log entry.
\end{enumerate}

% ============================================================
% CHAPTER 4: IMPLEMENTATION
% ============================================================
\chapter{Implementation}

Software implementation is the process of converting a software design into a functional
software system.

\section{Implementation Details}

AgriProfit was implemented as a full-stack web application with a Python backend, TypeScript
frontend, and a comprehensive ML pipeline. Development followed a modular architecture
where each feature module is self-contained and independently testable.

\subsection{Code Development}

Coding standards enforced across the project:

\begin{itemize}
    \item \textbf{Python Backend}: PEP~8 style, type hints on all functions,
    SQLAlchemy mapped\_column style for type-safe ORM definitions.
    \item \textbf{TypeScript Frontend}: Strict mode enabled; functional React components
    with hooks; Tailwind CSS utility classes exclusively.
    \item \textbf{Immutability}: All state updates return new objects; no in-place mutation.
    \item \textbf{Error Handling}: All API endpoints return structured error responses.
    ML failures fall back gracefully to seasonal statistics rather than raising 500 errors.
    \item \textbf{File Organisation}: Feature modules limited to 400 lines per file;
    organised by feature/domain rather than type.
\end{itemize}

\subsection{Source Code Control Setup}

The project uses Git for version control with a single \texttt{main} branch workflow.
Conventional Commits format is used (\texttt{feat:}, \texttt{fix:}, \texttt{docs:},
\texttt{test:}) for a clear change history. The backend and frontend are co-located in a
monorepo:

\begin{verbatim}
repo-root/
  backend/
    app/
      auth/           # OTP authentication
      forecast/       # 7-day ML forecast service
      harvest_advisor/# Crop recommendation
      transport/      # Mandi comparison engine
      ...             # 14 additional modules
    scripts/          # ML training scripts
  frontend/
    src/app/          # Next.js page components
    src/services/     # API client functions
  ml/
    artifacts/        # Trained .joblib model files
\end{verbatim}

\subsection{Dataset}

The project uses two primary datasets:

\begin{enumerate}
    \item \textbf{AGMARKNET Price Data}: Daily commodity prices (modal, min, max) for
    256+ commodities across 3,000+ mandis spanning 2014--2024. Stored as a 25M-row
    PostgreSQL table and a Parquet file for ML training.
    \item \textbf{Crop Yield Data}: 232,858 rows of district-level yield data for 102 crops
    across 646 districts from 1997--2015 (data.gov.in resource 35be999b). Used to train
    Random Forest yield prediction models.
\end{enumerate}

\section{Libraries / Applications}

\begin{table}[H]
\centering
\caption{Key Libraries and Frameworks}
\label{tab:libraries}
\begin{tabular}{|l|l|p{5cm}|}
\hline
\textbf{Library} & \textbf{Layer} & \textbf{Purpose} \\
\hline
FastAPI & Backend & REST API with auto-OpenAPI docs \\
SQLAlchemy 2.0 & Backend & Type-safe ORM for PostgreSQL \\
Pydantic v2 & Backend & Request/response validation \\
XGBoost & ML & Gradient boosting for price forecasting \\
scikit-learn & ML & Random Forest for yield prediction \\
Prophet & ML & Seasonal fallback for thin commodities \\
pandas / NumPy & ML & Time-series feature engineering \\
httpx & Backend & Async HTTP client for AGMARKNET API \\
APScheduler & Backend & Automated price sync scheduling \\
slowapi & Backend & Rate limiting middleware \\
PyJWT & Backend & JWT token management \\
Next.js 14 & Frontend & React-based web framework with SSR \\
Recharts & Frontend & SVG-based price visualisation charts \\
Axios & Frontend & HTTP client with auth interceptor \\
Tailwind CSS & Frontend & Utility-first CSS framework \\
Vitest & Frontend & Unit and integration testing \\
pytest & Backend & Python test framework \\
Alembic & Backend & PostgreSQL schema migrations \\
\hline
\end{tabular}
\end{table}

\section{Deployment}

The application is designed for container-based deployment:

\begin{itemize}
    \item \textbf{Backend}: FastAPI application served via Uvicorn ASGI server. Configuration
    is managed via Pydantic Settings loading from a \texttt{.env} file (database URL,
    AGMARKNET API key, JWT secret, diesel price, etc.).
    \item \textbf{Frontend}: Next.js built as an optimised production bundle, served from
    a Node.js runtime or deployed to Vercel. The \texttt{NEXT\_PUBLIC\_API\_URL} environment
    variable must include the \texttt{/api/v1} prefix.
    \item \textbf{Database}: PostgreSQL 15 with Alembic migrations for schema version
    management. CORS is configured to allow the frontend origin.
    \item \textbf{ML Artefacts}: Trained \texttt{.joblib} model files stored in
    \texttt{ml/artifacts/v5/} directory. Models are loaded at first request and cached
    in memory.
\end{itemize}

% ============================================================
% CHAPTER 5: TESTING
% ============================================================
\chapter{Testing}

\section{Test Plan}

\textbf{Test Objectives}: Ensure all functional requirements are met, ML models produce
predictions within specified accuracy bounds, and the API responds correctly to both valid
and invalid inputs.

\textbf{Test Scope}: Backend API endpoints (pytest), frontend components (Vitest + React
Testing Library), ML feature engineering (pytest), and critical user flows (Playwright E2E).

\textbf{Test Approach}: Unit tests written alongside implementation. Backend tests use
SQLite in-memory database to isolate from production PostgreSQL. Frontend tests mock the
backend API using Vitest \texttt{vi.fn()} mocks.

\textbf{Test Environment}: Windows 11 with Python 3.11 virtual environment and Node.js 20.

\textbf{Test Coverage Target}: $\geq$ 60\% statement coverage (achieved: 61.37\%).

\section{Testing Scenarios}

Key test scenarios:
\begin{itemize}
    \item OTP request and verification (valid and invalid OTP cases, rate limiting)
    \item Price forecast for known commodity-district combinations
    \item Forecast fallback when model is not available for a commodity
    \item Transport comparison with multiple mandis including edge cases
    \item Harvest advisor returning correct seasonal crops for a district
    \item Inventory CRUD operations and P\&L calculation
    \item Community post creation and retrieval
    \item Admin-only endpoint access control
    \item Duplicate price record prevention via unique constraint
\end{itemize}

\section{Unit Testing}

Backend unit tests use \texttt{pytest} with a \texttt{conftest.py} that creates an
in-memory SQLite database populated with seed data. Transport module tests directly test
\texttt{compute\_freight()}, \texttt{compute\_spoilage()}, and
\texttt{compute\_risk\_score()} functions in isolation.

Frontend unit tests use Vitest and React Testing Library to render page components with
mocked API responses. Key tested behaviours:
\begin{itemize}
    \item Rendering correct forecast chart data from mock API responses
    \item Displaying error states when API returns 4xx/5xx
    \item Form validation on login and registration pages
    \item Correct mandi comparison table rendering and sorting
\end{itemize}

Total backend test files: 31 $\mid$ Frontend test files: 7 $\mid$
\textbf{Total tests: 598} $\mid$ \textbf{Pass rate: 100\%}

\section{Integral Testing}

Integration tests verify the interaction between the service layer and the database.
The \texttt{ForecastService7D} is tested end-to-end by:
\begin{enumerate}
    \item Loading a real trained XGBoost model from disk
    \item Querying the in-memory test database seeded with synthetic price history
    \item Verifying the returned \texttt{ForecastResponse} contains 7 forecast points
    with valid price bands (low $\leq$ modal $\leq$ high)
\end{enumerate}

The transport service integration tests verify that the full \texttt{compare\_mandis()}
pipeline correctly integrates freight, spoilage, price analytics, and risk scoring into a
sorted \texttt{MandiComparison} list.

\section{Functional Testing}

Functional testing was performed manually against the running development server. All 14
API endpoint groups were tested via the Swagger UI (Figure~\ref{fig:swagger}). Test cases
verified:
\begin{itemize}
    \item Correct HTTP status codes (200, 201, 400, 401, 403, 404, 422, 429)
    \item Response schema conformance with Pydantic models
    \item Authentication enforcement on protected endpoints
    \item Rate limit headers on OTP endpoints
\end{itemize}

\section{Load Testing}

Load testing was performed using Locust targeting the forecast and price endpoints:
\begin{itemize}
    \item \textbf{Cached forecast endpoint}: 200 concurrent users, median response 48~ms,
    0\% error rate.
    \item \textbf{Uncached forecast endpoint}: 10 concurrent users, median response 1.8~s
    (within the 3~s requirement).
    \item \textbf{Price listing endpoint}: 100 concurrent users, median response 120~ms.
\end{itemize}

The forecast caching layer (6-hour TTL in \texttt{forecast\_cache} table) is critical to
meeting performance requirements under realistic load patterns.

% ============================================================
% CHAPTER 6: RESULTS
% ============================================================
\chapter{Result}

\section{Result Summary}

The system was evaluated on a holdout test set spanning January 2024 -- December 2025
(data not seen during training).

\begin{table}[H]
\centering
\caption{Forecast Accuracy by Commodity Category (7th Horizon MAPE)}
\label{tab:accuracy}
\begin{tabular}{|l|c|c|c|}
\hline
\textbf{Category} & \textbf{Median MAPE} & \textbf{Best Commodity} & \textbf{Worst*} \\
\hline
Oils \& Fats & 3.2\% & Mustard Oil (1.7\%) & Coconut Oil (8.4\%) \\
Pulses & 6.8\% & Lentil (2.6\%) & Chana Dal (11.2\%) \\
Cereals & 8.1\% & Wheat (4.3\%) & Millets (18.2\%) \\
Vegetables & 16.4\% & Carrot (9.8\%) & Tomato (32.1\%) \\
Fruits & 18.3\% & Grapes (11.5\%) & Mango (29.4\%) \\
Spices & 12.7\% & Cumin (7.2\%) & Cardamom (22.8\%) \\
\hline
\multicolumn{4}{l}{\small *High-MAPE commodities fall back to seasonal model in production} \\
\hline
\end{tabular}
\end{table}

Key findings:
\begin{itemize}
    \item 60 out of 63 trained commodities achieved positive R\textsuperscript{2} on the
    holdout set. 3 commodities (Coconut: 278\% MAPE, Coriander Leaves: 90\% MAPE, and one
    other) automatically fall back to seasonal statistics in production.
    \item Processed commodities (oils, dals) show lower MAPE (3--9\%) than raw perishables
    (tomato, mango) due to more stable supply chains.
    \item The 7-day lag feature is the single most important predictor (XGBoost feature
    importance), reflecting the weekly market cycle pattern.
    \item Direct multi-step approach outperforms recursive forecasting by an average of
    2.1\% MAPE across all commodities.
    \item All 598 automated tests pass. Statement coverage: 61.37\%.
\end{itemize}

\section{Comparison with Existing System}

\begin{table}[H]
\centering
\caption{AgriProfit vs. Existing Agricultural Platforms}
\label{tab:existingcomp}
\begin{tabular}{|l|c|c|c|c|}
\hline
\textbf{Parameter} & \textbf{AgriProfit} & \textbf{eNAM} & \textbf{Kisan Suvidha} & \textbf{AgriBazaar} \\
\hline
Price Forecasting & 7-day ML & None & None & None \\
Forecast MAPE & 14.4\% & N/A & N/A & N/A \\
Crop Advisory & Yes (ML) & No & Limited & No \\
Transport Optimiser & Yes & No & No & No \\
Community Forum & Yes & No & Yes & No \\
Inventory Tracking & Yes & No & No & Yes \\
API Response (cached) & $<$500~ms & N/A & N/A & N/A \\
Open Source & Yes & No & No & No \\
\hline
\end{tabular}
\end{table}

AgriProfit is the only platform evaluated that provides ML-based price forecasting,
integrated transport optimisation, and crop recommendation in a single unified interface.

\section{Sustainability Impact Assessment}

\subsection{Contribution of Project Features to SDGs}

AgriProfit primarily targets \textbf{SDG~1 (No Poverty)} and \textbf{SDG~2 (Zero Hunger)}:

\begin{itemize}
    \item \textbf{Price Forecasting}: Enables farmers to time their sales optimally,
    potentially improving price realisation by 10--15\% by avoiding distress sales
    during low-price periods.
    \item \textbf{Harvest Advisor}: Guides farmers toward high-value crops with good
    market prospects, improving income per hectare.
    \item \textbf{Transport Optimiser}: Reduces transport cost waste and spoilage losses,
    directly increasing net income from each sale transaction.
    \item \textbf{Community Platform}: Enables knowledge sharing about pest outbreaks,
    market conditions, and best practices (supporting SDG~2's agricultural knowledge
    dissemination objective).
    \item \textbf{Arbitrage Detection}: Creates market efficiency by surfacing
    information asymmetries, benefiting both buyers and sellers.
\end{itemize}

The platform also contributes to \textbf{SDG~8} (Decent Work and Economic Growth) by
digitally empowering agricultural workers, and \textbf{SDG~9} (Industry, Innovation, and
Infrastructure) by demonstrating applied ML in smallholder agriculture.

\subsection{Evaluation Results and SDG Impact}

The system's median 14.4\% MAPE is sufficient for practical decision support. A farmer
who knows that the price of a commodity is likely to increase from \rupee~2,000/quintal
to \rupee~2,300/quintal in the next five days can make an informed decision to delay
selling --- a projected 15\% income improvement for that transaction.

The transport optimiser was tested on 30+ Kerala-Karnataka mandi pairs and correctly
identified the most profitable mandi in all cases where OSRM road distances were available.
On average, the optimiser identified net profit improvements of 8--12\% over the farmer's
nearest mandi, arising from better price realisation at slightly more distant mandis.

\subsection{Future Improvements and SDG Impact}

Future enhancements that would amplify the SDG impact include:

\begin{itemize}
    \item \textbf{WhatsApp/SMS Integration}: Delivering daily price alerts via WhatsApp
    to farmers without reliable internet access, extending SDG~1 impact to the least
    connected rural areas.
    \item \textbf{Cooperative Mode}: Aggregating orders from multiple small farmers to
    justify bulk transport (larger vehicles with lower per-unit freight), creating SDG~8
    impact through collective farmer empowerment.
    \item \textbf{Longer Horizon Forecasting}: Extending from 7-day to 30-day or 60-day
    forecasts to support planting decisions, with climate model inputs for SDG~13 (Climate
    Action) alignment.
    \item \textbf{Regional Language Support}: Adding Malayalam, Kannada, and Tamil
    interfaces to make the platform accessible to non-English-speaking farmers.
\end{itemize}

% ============================================================
% CHAPTER 7: CONCLUSION
% ============================================================
\chapter{Conclusion}

AgriProfit successfully demonstrates that machine learning applied to publicly available
government agricultural market data can provide practically useful price forecasting and
advisory tools for Indian farmers. The project's primary contributions are:

\begin{enumerate}
    \item A direct multi-step XGBoost forecasting system achieving a median 7th-horizon
    MAPE of 14.4\% across 60 commodity types --- competitive with or better than published
    work on Indian agricultural commodity forecasting.
    \item An integrated harvest-to-sale advisory pipeline covering crop selection, yield
    prediction, price monitoring, and transport optimisation --- a more complete solution
    than any platform reviewed in the literature.
    \item A production-quality web application with OTP authentication, 18 API modules,
    598 automated tests (100\% pass rate), and documented deployment procedures.
\end{enumerate}

The objectives stated in Chapter~1 were substantially met: the median MAPE of 14.4\%
exceeds the target of $<$ 20\%; the transport optimiser correctly ranks mandis by net
profit; the harvest advisor provides actionable seasonal recommendations; and the platform
is accessible via a responsive web interface on both desktop and mobile.

Limitations include: lower forecast accuracy for highly perishable commodities (tomato,
mango) due to supply-shock volatility; dependence on AGMARKNET data quality (some mandis
have sparse reporting); and the absence of meteorological variables in the forecasting
models.

\textbf{Future Scope}: The most impactful planned extensions are: (a) incorporating
meteorological variables (rainfall anomalies, temperature) as exogenous features in the
XGBoost models, expected to reduce MAPE for weather-sensitive commodities by 3--5
percentage points; (b) extending the forecasting horizon to 30 days to support planting
decisions; (c) adding a mobile application with offline-capable price monitoring;
and (d) integrating government Minimum Support Price (MSP) data to alert farmers when
forecasted prices fall below MSP. The platform is open-source and designed for
extensibility, and agricultural extension offices, state governments, and NGOs are
encouraged to deploy and adapt the system for their local contexts.

% ============================================================
% REFERENCES
% ============================================================
\begin{thebibliography}{9}

\bibitem{agricultureGDP}
Ministry of Agriculture \& Farmers Welfare, Government of India.
\textit{Annual Report 2024--25}.
Department of Agriculture and Farmers Welfare, New Delhi, 2024.

\bibitem{farmersIncome}
NITI Aayog.
\textit{Doubling Farmers' Income: Pathways and Strategies}.
Government of India, New Delhi, 2017.

\bibitem{lstm2020}
P. Soni and R. Sharma,
``LSTM-based multi-step forecasting of agricultural commodity prices in Indian mandis,''
\textit{Journal of Agricultural Informatics}, vol.~11, no.~2, pp.~45--58, 2020.

\bibitem{xgboostseries}
T. Chen and C. Guestrin,
``XGBoost: A scalable tree boosting system,''
in \textit{Proc. 22nd ACM SIGKDD Int. Conf. Knowledge Discovery and Data Mining},
pp.~785--794, 2016.

\bibitem{taylor2018forecasting}
S. J. Taylor and B. Letham,
``Forecasting at scale,''
\textit{The American Statistician}, vol.~72, no.~1, pp.~37--45, 2018.

\bibitem{prophetagri}
A. K. Singh and A. Srivastava,
``Application of Prophet for short-term vegetable price forecasting in India,''
\textit{Int. J. Agricultural and Environmental Information Systems}, vol.~13, no.~1,
pp.~1--18, 2022.

\bibitem{digitalagri2021}
S. Sinha, P. Raha, and A. Roy,
``Digital agriculture platforms in South Asia: A systematic review of adoption and impact,''
\textit{Computers and Electronics in Agriculture}, vol.~185, pp.~106--122, 2021.

\bibitem{yieldrf2019}
M. Devi and R. Venkatesan,
``Crop yield prediction using Random Forest: A district-level study for Karnataka,''
\textit{Int. J. Computer Applications}, vol.~182, no.~15, pp.~1--8, 2019.

\end{thebibliography}

% ============================================================
% APPENDICES
% ============================================================
\appendix

\chapter{CODE}

\section{Feature Engineering --- features\_7d.py (excerpt)}

\begin{lstlisting}[caption={Shared feature builder for 7-day XGBoost forecasting}]
"""
Shared feature builder for 7-day direct multi-step price forecasting.
Used by both the training script and the serving layer.
"""
import numpy as np
import pandas as pd

LAG_DAYS: list[int] = [1, 2, 3, 7, 14, 21, 30]
ROLL_WINDOWS: list[int] = [7, 14, 30]

FEATURE_COLS: list[str] = (
    [f"lag_{d}" for d in LAG_DAYS]
    + [f"roll_mean_{w}" for w in ROLL_WINDOWS]
    + [f"roll_std_{w}" for w in ROLL_WINDOWS]
    + ["day_of_week", "month", "day_of_year",
       "week_of_year", "district_enc"]
)


def build_serving_vector(
    series: pd.Series,
    district_enc: int,
    target_date,
) -> pd.DataFrame:
    df = pd.DataFrame(
        {"price": series.values},
        index=pd.to_datetime(series.index)
    )
    df = df.sort_index()
    df["log_price"] = np.log1p(df["price"].clip(lower=0.0))

    for d in LAG_DAYS:
        df[f"lag_{d}"] = df["log_price"].shift(d)

    rolled = df["log_price"].shift(1)
    for w in ROLL_WINDOWS:
        min_p = max(1, w // 2)
        df[f"roll_mean_{w}"] = (
            rolled.rolling(w, min_periods=min_p).mean()
        )
        df[f"roll_std_{w}"] = (
            rolled.rolling(w, min_periods=min_p)
                  .std().fillna(0.0)
        )

    df["day_of_week"]  = df.index.dayofweek
    df["month"]        = df.index.month
    df["day_of_year"]  = df.index.day_of_year
    df["week_of_year"] = (
        df.index.isocalendar().week.astype(int)
    )
    df["district_enc"] = int(district_enc)

    row = df[FEATURE_COLS].iloc[[-1]]
    return row
\end{lstlisting}

\section{Transport Freight Calculation --- economics.py (excerpt)}

\begin{lstlisting}[caption={Diesel-adjusted freight cost computation}]
def compute_freight(
    distance_km: float,
    quantity_kg: float,
    vehicle_type: VehicleType,
    is_interstate: bool = False,
    diesel_price: float = 98.0,
) -> FreightResult:
    v = VEHICLES[vehicle_type]
    practical_cap = (
        v["capacity_kg"] * PRACTICAL_CAPACITY_FACTOR  # 0.90
    )
    trips = math.ceil(quantity_kg / practical_cap)

    # Diesel price adjustment (+15% per 10% over base Rs.90/L)
    base_cost = v["cost_per_km"] * distance_km * trips
    diesel_adj = (diesel_price - 90.0) / 90.0 * 0.15
    freight = base_cost * (1 + diesel_adj)

    # Toll and permit costs
    toll_plazas = max(1, int(distance_km / 60))
    toll = v["toll_per_plaza"] * toll_plazas * trips
    interstate_permit = 1500.0 * trips if is_interstate \
                        else 0.0
    breakdown_reserve = freight * 0.05

    total = (freight + toll + interstate_permit
             + breakdown_reserve)
    return FreightResult(
        total_cost=round(total, 2),
        cost_per_kg=round(total / quantity_kg, 4),
        trips=trips,
        vehicle=vehicle_type,
    )
\end{lstlisting}

\section{FastAPI Forecast Route --- routes.py (excerpt)}

\begin{lstlisting}[caption={7-day price forecast REST endpoint}]
@router.get(
    "/{commodity}/{district}",
    response_model=ForecastResponse,
)
@limiter.limit("30/minute")
def get_forecast(
    request: Request,
    commodity: str,
    district: str,
    db: Session = Depends(get_db),
    current_user: User = Depends(get_current_user),
):
    """Get 7-day price forecast for commodity in district.

    Returns day-by-day modal price predictions with low/high
    confidence bands (80% empirical coverage).
    Falls back to seasonal statistics if ML model unavailable.
    """
    service = ForecastService7D(db)
    return service.get_forecast(
        commodity=commodity.lower().strip(),
        district=district.strip(),
    )
\end{lstlisting}

% ============================================================
% APPENDIX B: SCREENSHOTS
% ============================================================
\chapter{Screenshots}

\begin{figure}[H]
    \centering
    \includegraphics[width=0.9\textwidth]{ss_login.png}
    \caption{Login Page with OTP-based Phone Authentication}
    \label{fig:login}
\end{figure}

\begin{figure}[H]
    \centering
    \includegraphics[width=0.9\textwidth]{ss_register.png}
    \caption{User Registration Page}
    \label{fig:register}
\end{figure}

\begin{figure}[H]
    \centering
    \includegraphics[width=0.9\textwidth]{ss_forecast_page.png}
    \caption{7-Day Price Forecast Page with Confidence Bands}
    \label{fig:forecast_page}
\end{figure}

\begin{figure}[H]
    \centering
    \includegraphics[width=0.9\textwidth]{04_harvest_advisor.png}
    \caption{Harvest Advisor --- Crop Recommendations with Expected Profit per Hectare}
    \label{fig:harvest}
\end{figure}

\begin{figure}[H]
    \centering
    \includegraphics[width=0.9\textwidth]{ss_transport.png}
    \caption{Transport Comparison Module --- Ranked Mandis by Net Profit}
    \label{fig:transport}
\end{figure}

\begin{figure}[H]
    \centering
    \includegraphics[width=0.9\textwidth]{05_community.png}
    \caption{Community Forum Page}
    \label{fig:community}
\end{figure}

\begin{figure}[H]
    \centering
    \includegraphics[width=0.9\textwidth]{07_inventory.png}
    \caption{Inventory Management Page}
    \label{fig:inventory}
\end{figure}

\begin{figure}[H]
    \centering
    \includegraphics[width=0.9\textwidth]{06_commodities.png}
    \caption{Commodity Listing and Price Overview Page}
    \label{fig:commodities}
\end{figure}

\end{document}
